\documentclass[a4paper, 14pt]{extarticle}

\usepackage{cmap}
\usepackage[T2A]{fontenc}
\usepackage[utf8]{inputenc}
\usepackage[english, russian]{babel}
\usepackage{amsmath}
\usepackage{subcaption}
\usepackage{graphicx}
\usepackage{float}
\usepackage{hyperref}

\usepackage[left=3cm, right=1.5cm, top=2cm, bottom=2cm]{geometry}
\usepackage{setspace}
\onehalfspacing 

\usepackage{indentfirst}
\setlength{\parindent}{1.25cm} 
\setlength{\emergencystretch}{3em}
\usepackage[expansion=false,protrusion=false,nopatch=footnote]{microtype}

\usepackage{tocloft}
\addto\captionsrussian{\renewcommand{\contentsname}{СОДЕРЖАНИЕ}}
\renewcommand{\cfttoctitlefont}{\hfill\normalfont}
\renewcommand{\cftaftertoctitle}{\hfill\vspace{\baselineskip}}
\renewcommand{\cftsecdotsep}{\cftdotsep}

\begin{document}

\begin{titlepage}
    \begin{center}
        Министерство науки и высшего образования Российской Федерации\\
        Федеральное государственное автономное образовательное учреждение высшего образования\\
        «Санкт-Петербургский политехнический университет Петра Великого»\\
        
        \textbf{Физико-механический институт}\\
        \textbf{Высшая школа прикладной математики и вычислительной физики}
    \end{center}

    \vfill
    
    \begin{center}
        \textbf{Отчёт по лабораторной работе \textnumero~ 2\\[0.5cm]}
        \textbf{Характеристики положения и рассеяния\\}
        по дисциплине <<Математическая статистика>>
    \end{center}

    \vfill

    \hfill
    \begin{minipage}{0.45\textwidth}
        \textbf{Выполнил студент гр. 5030102/30202:}\\
        Шильдт Д.С.\\[0.5cm]
        \textbf{Преподаватель:}\\ 
        Баженов А.Н.
    \end{minipage}

    \vfill

    \begin{center}
        Санкт-Петербург\\2026
    \end{center}
\end{titlepage}

\setcounter{page}{2}

\tableofcontents
\newpage

\section{Постановка задачи}

\textbf{Цели и задачи работы:}
\begin{itemize}
    \item Исследовать сходимость выборочных характеристик к теоретическим при росте $n$.
    \item Исследовать оценки характеристик положения на устойчивость к выбросам.
\end{itemize}

\textbf{Задание:}
Сгенерировать выборки объёмом $n = 10, 100, 1000$. Для каждой выборки вычислить:
\begin{enumerate}
    \item Выборочное среднее $\bar{x}$;
    \item Выборочную медиану $\text{med}$;
    \item Полусумму экстремальных элементов $z_R = \frac{x_{\min} + x_{\max}}{2}$;
    \item Полусумму квартилей $z_Q = \frac{z_{1/4} + z_{3/4}}{2}$;
    \item Усеченное среднее $z_{tr}$ (с отбрасыванием $10\%$ наименьших и наибольших значений).
\end{enumerate}

Повторить генерацию и вычисления 1000 раз. Найти среднее и дисперсию каждой характеристики: $$E(z) = \overline{z},\quad D(z) = \overline{z^2} - \overline{z}^2$$

\section{Теоретическая часть}

Основой для вычисления точечных оценок характеристик положения служит вариационный ряд, представляющий собой элементы выборки, упорядоченные по неубыванию: $x_{(1)} \le x_{(2)} \le \dots \le x_{(n)}$.

В работе исследуются пять статистических оценок истинного центра распределения $m$ \cite{babakhina, maksimov}:

1. \textbf{Выборочное среднее:}
\begin{equation}
    \bar{x} = \frac{1}{n}\sum_{i=1}^n x_i
\end{equation}

2. \textbf{Выборочная медиана:}
\begin{equation}
    med = 
    \begin{cases}
        x_{(l+1)}, & \text{при } n = 2l + 1 \\
        \frac{x_{(l)} + x_{(l+1)}}{2}, & \text{при } n = 2l
    \end{cases}
\end{equation}

3. \textbf{Полусумма экстремальных выборочных элементов:}
\begin{equation}
    z_R = \frac{x_{(1)} + x_{(n)}}{2}
\end{equation}

4. \textbf{Полусумма квартилей:}
\begin{equation}
    z_Q = \frac{z_{1/4} + z_{3/4}}{2}
\end{equation}
где $z_p$ — выборочная квантиль порядка $p$.

5. \textbf{Усеченное среднее:}
\begin{equation}
    z_{tr} = \frac{1}{n-2r} \sum_{i=r+1}^{n-r} x_{(i)} \qquad r\approx np
\end{equation}
где $r = [n \cdot p]$. В расчетах применяется симметричное отбрасывание $10\%$ элементов с каждой стороны ($p=0.1$) [1].

\section{Реализация}

\subsection{Язык программирования и среда}
    \begin{itemize}
        \item \textbf{Используемый язык:} Python 3.14.2
        \item \textbf{Используемая среда разработки:} Visual Studio Code
    \end{itemize}

\subsection{Используемые библиотеки}
    \begin{itemize}
        \item \textbf{numpy} --- использовалась для работы с массивами данных и математическими функциями.
        \item \textbf{scipy} --- использовалась для генерации выборок из теоретических распределений и вычисления функций плотности/вероятности.
        \item \textbf{matplotlib} --- использовалась для визуализации данных и построения гистограмм.
        \item \textbf{pathlib} --- использовалась для автоматического создания каталогов и сохранения итоговых графиков.
    \end{itemize}

\subsection{Суть реализованных методов}

\begin{itemize}
    \item \texttt{calculateCharacteristics(sample)} — вычисляет вектор из пяти точечных оценок положения ($\bar{x}$, $med$, $z_R$, $z_Q$, $z_{tr}$) для переданной выборки. Предварительно массив сортируется по неубыванию для формирования вариационного ряда. 
    \item \texttt{estimateMoments(values)} — оценивает эмпирическое математическое ожидание $E(z)$ и выборочную дисперсию $D(z)$ по множеству реализаций каждой оценки. Дисперсия вычисляется через разность моментов: $D(z) = E(z^2) - (E(z))^2$.
    \item \texttt{runExperiment(dist\_name, generator, rng)} — реализует логику многократного статистического моделирования. Для каждого объема выборки $n \in \{10, 100, 1000\}$ инициирует внутренний цикл из $1000$ итераций генерации данных. Накопленные массивы характеристик затем передаются в \texttt{estimateMoments} для финальной агрегации.
    \item \texttt{saveCsv(rows)} и \texttt{printReport(rows)} — отвечают за форматирование и экспорт итоговых значений.
\end{itemize}

\section{Результаты}

В данном разделе представлены сводные результаты вычисления точечных оценок характеристик положения. Данные получены путем многократного статистического моделирования (1000 повторений генерации) для каждого объема выборки $n \in \{10, 100, 1000\}$. 

Итоговые значения в таблицах \ref{tab:res_norm}--\ref{tab:res_uniform} записаны в формате $E \pm \sqrt{D}$, где $E$ — эмпирическое математическое ожидание точечной оценки, $D$ — выборочная дисперсия.

\begin{table}[H]
    \centering
    \caption{Характеристики положения для нормального распределения $N(0, 1)$}
    \label{tab:res_norm}
    \begin{tabular}{|l|c|c|c|}
        \hline
        \textbf{Характеристика} & $\mathbf{n=10}$ & $\mathbf{n=100}$ & $\mathbf{n=1000}$ \\ \hline
        Выборочное среднее $\bar{x}$ & $0.0 \pm 0.3$ & $0.0 \pm 0.1$ & $0.0 \pm 0.0$ \\ \hline
        Медиана $med$ & $0.0 \pm 0.4$ & $0.0 \pm 0.1$ & $0.0 \pm 0.0$ \\ \hline
        $z_R$ & $0.0 \pm 0.4$ & $0.0 \pm 0.3$ & $0.0 \pm 0.3$ \\ \hline
        $z_Q$ & $0.0 \pm 0.3$ & $0.0 \pm 0.1$ & $0.0 \pm 0.0$ \\ \hline
        $z_{tr}$ & $0.0 \pm 0.3$ & $0.0 \pm 0.1$ & $0.0 \pm 0.0$ \\ \hline
    \end{tabular}
\end{table}

\begin{table}[H]
    \centering
    \caption{Характеристики положения для распределения Коши $C(0, 1)$}
    \label{tab:res_cauchy}
    \begin{tabular}{|l|c|c|c|}
        \hline
        \textbf{Характеристика} & $\mathbf{n=10}$ & $\mathbf{n=100}$ & $\mathbf{n=1000}$ \\ \hline
        Выборочное среднее $\bar{x}$ & $0.5 \pm 51.5$ & $1.5 \pm 50.5$ & $5.4 \pm 131.2$ \\ \hline
        Медиана $med$ & $0.0 \pm 0.6$ & $0.0 \pm 0.2$ & $0.0 \pm 0.0$ \\ \hline
        $z_R$ & $2.4 \pm 257.0$ & $73.6 \pm 2521.5$ & $2678.2 \pm 65532.6$ \\ \hline
        $z_Q$ & $0.0 \pm 1.0$ & $0.0 \pm 0.2$ & $0.0 \pm 0.1$ \\ \hline
        $z_{tr}$ & $0.0 \pm 1.3$ & $0.0 \pm 0.2$ & $0.0 \pm 0.1$ \\ \hline
    \end{tabular}
\end{table}

\begin{table}[H]
    \centering
    \caption{Характеристики положения для распределения Лапласа $L(0, 1/\sqrt{2})$}
    \label{tab:res_laplace}
    \begin{tabular}{|l|c|c|c|}
        \hline
        \textbf{Характеристика} & $\mathbf{n=10}$ & $\mathbf{n=100}$ & $\mathbf{n=1000}$ \\ \hline
        Выборочное среднее $\bar{x}$ & $0.0 \pm 0.3$ & $0.0 \pm 0.1$ & $0.0 \pm 0.0$ \\ \hline
        Медиана $med$ & $0.0 \pm 0.3$ & $0.0 \pm 0.1$ & $0.0 \pm 0.0$ \\ \hline
        $z_R$ & $0.0 \pm 0.7$ & $0.0 \pm 0.6$ & $0.0 \pm 0.6$ \\ \hline
        $z_Q$ & $0.0 \pm 0.3$ & $0.0 \pm 0.1$ & $0.0 \pm 0.0$ \\ \hline
        $z_{tr}$ & $0.0 \pm 0.3$ & $0.0 \pm 0.1$ & $0.0 \pm 0.0$ \\ \hline
    \end{tabular}
\end{table}

\begin{table}[H]
    \centering
    \caption{Характеристики положения для распределения Пуассона $P(5)$}
    \label{tab:res_poisson}
    \begin{tabular}{|l|c|c|c|}
        \hline
        \textbf{Характеристика} & $\mathbf{n=10}$ & $\mathbf{n=100}$ & $\mathbf{n=1000}$ \\ \hline
        Выборочное среднее $\bar{x}$ & $5.0 \pm 0.7$ & $5.0 \pm 0.2$ & $5.0 \pm 0.1$ \\ \hline
        Медиана $med$ & $4.9 \pm 0.8$ & $4.9 \pm 0.3$ & $5.0 \pm 0.0$ \\ \hline
        $z_R$ & $5.3 \pm 1.0$ & $6.0 \pm 0.7$ & $6.8 \pm 0.6$ \\ \hline
        $z_Q$ & $4.9 \pm 0.7$ & $4.9 \pm 0.4$ & $4.6 \pm 0.3$ \\ \hline
        $z_{tr}$ & $4.9 \pm 0.7$ & $4.9 \pm 0.2$ & $4.9 \pm 0.1$ \\ \hline
    \end{tabular}
\end{table}

\begin{table}[H]
    \centering
    \caption{Характеристики положения для равномерного распределения $U(-\sqrt{3}, \sqrt{3})$}
    \label{tab:res_uniform}
    \begin{tabular}{|l|c|c|c|}
        \hline
        \textbf{Характеристика} & $\mathbf{n=10}$ & $\mathbf{n=100}$ & $\mathbf{n=1000}$ \\ \hline
        Выборочное среднее $\bar{x}$ & $0.0 \pm 0.3$ & $0.0 \pm 0.1$ & $0.0 \pm 0.0$ \\ \hline
        Медиана $med$ & $0.0 \pm 0.5$ & $0.0 \pm 0.2$ & $0.0 \pm 0.1$ \\ \hline
        $z_R$ & $0.0 \pm 0.2$ & $0.0 \pm 0.0$ & $0.0 \pm 0.0$ \\ \hline
        $z_Q$ & $0.0 \pm 0.4$ & $0.0 \pm 0.1$ & $0.0 \pm 0.0$ \\ \hline
        $z_{tr}$ & $0.0 \pm 0.3$ & $0.0 \pm 0.1$ & $0.0 \pm 0.0$ \\ \hline
    \end{tabular}
\end{table}
\section{Обсуждение}

Анализ сводных результатов многократного статистического моделирования позволяет оценить состоятельность, эффективность и робастность исследуемых точечных характеристик положения.

Для распределений, обладающих конечными теоретическими моментами (Нормальное, Лапласа, Равномерное, Пуассона), эмпирические математические ожидания $E(z)$ всех пяти оценок сгруппированы вблизи истинных генеральных центров ($m=0$ и $m=5$ соответственно). Сходимость выборочных оценок к заданным параметрам подтверждает их состоятельность.

Теоретическая закономерность убывания дисперсии выборочного среднего $D(\bar{x}) = \frac{\sigma^2}{n}$ [1] строго прослеживается на полученных эмпирических данных. Для стандартного нормального распределения $N(0, 1)$ расчетная дисперсия $\bar{x}$ последовательно снижается обратно пропорционально росту $n$: от $0.0997$ при $n=10$ до $0.0102$ при $n=100$, достигая $0.000898$ при $n=1000$. Полученное при $n=1000$ значение практически совпадает с теоретическим асимптотическим пределом $0.001$. Аналогичная сходимость дисперсий наблюдается у всех характеристик на симметричных распределениях без выраженных аномалий.

Критическое различие между статистическими оценками выявляется при анализе их робастности — устойчивости к аномальным значениям [2]. Выборочное среднее $\bar{x}$ и полусумма экстремальных элементов $z_R$ не защищены от влияния выбросов [3]. Данный факт наглядно демонстрируют расчеты для распределения Коши $C(0, 1)$. Генерируемые законом Коши «тяжелые хвосты» приводят к неограниченному росту дисперсий неробастных оценок. При $n=1000$ расчетная дисперсия выборочного среднего достигает $17221.5$, а дисперсия $z_R$ принимает аномальное значение $4.29 \cdot 10^9$. В таких условиях оценки полностью теряют статистический смысл.

В противовес этому, характеристики $med$, $z_Q$ и $z_{tr}$, базирующиеся на центральных порядковых статистиках вариационного ряда, демонстрируют высокую робастность [3]. Для того же распределения Коши при $n=1000$ их дисперсии сохраняют устойчивость и остаются минимальными: $D(med) = 0.0024$, $D(z_{tr}) = 0.0045$. 

Полученные цифры доказывают, что при наличии аномальных наблюдений или неизвестном характере "хвостов" распределения использование среднего арифметического $\bar{x}$ некорректно [2]. Для получения надежных результатов необходимо применять робастные оценки: выборочную медиану, полусумму квартилей или усеченное среднее.

\section{Выводы}

В результате многократного статистического моделирования исследованы свойства пяти точечных оценок характеристик положения. 

Практически подтверждено, что для распределений с конечными теоретическими моментами рассмотренные оценки состоятельны: при увеличении объема выборки $n$ их эмпирические математические ожидания группируются строго вокруг истинных генеральных центров, а выборочные дисперсии закономерно убывают.

Расчеты на базе распределения Коши доказали, что классическое выборочное среднее $\bar{x}$ и полусумма экстремальных элементов $z_R$ абсолютно не защищены от аномальных выбросов, из-за чего их применение для распределений с тяжелыми хвостами теряет статистический смысл. Установлено, что надежное оценивание центра группировки данных в условиях возможных аномалий обеспечивают только робастные характеристики, базирующиеся на центральных порядковых статистиках вариационного ряда: выборочная медиана $med$, полусумма квартилей $z_Q$ и усеченное среднее $z_{tr}$. 

 
\refstepcounter{section}
\addcontentsline{toc}{section}{\thesection\ Список литературы}
\renewcommand{\refname}{\thesection\ Список литературы}

\begin{thebibliography}{9}

    \bibitem{babakhina}
    Бабахина С. А., Баженов А. Н. Практикум по математической статистике: учеб. пособие. — СПб. : Санкт-Петербургский политехнический университет Петра Великого, 2026. 

    \bibitem{maksimov}
    Максимов Ю. Д. Математика. Теория и практика по математической статистике. Конспект-справочник по теории вероятностей : учеб. пособие. — СПб. : Изд-во Политехн. ун-та, 2009. — 395 с. 

\end{thebibliography}

\section{Приложение}

Исходный код программ доступен в репозитории на GitHub: \\
\url{https://github.com/Reichenau/math-stats-labs}
 
\end{document}
